\chapter{Introduction}

\section{Contexte}
Ce travail de Bachelor porte sur la conception complète d'un lecteur audio portable. Il couvre toute la chaîne de réalisation, du schéma électronique au routage de la carte, puis au firmware embarqué et enfin aux tests de validation. Le sujet a été choisi librement. Il ne répond pas à une commande industrielle, mais à la volonté de mener un projet d'ingénierie de bout en bout et d'obtenir un appareil réel, fabriqué et fonctionnel.

Le smartphone a remplacé le lecteur audio dédié sur le marché grand public. Ce type d'appareil subsiste toutefois dans des usages de niche, pour trois raisons. Il ne sert qu'à écouter de la musique, il n'apporte aucune distraction, et il fonctionne sans dépendre d'un autre équipement. L'appareil développé ici suit cette logique. Il est conçu comme un objet compact et autonome, et non comme une alternative au téléphone.


\section{Problématique}
Concevoir un lecteur audio portable ne se résume pas à assembler des blocs fonctionnels connus. La difficulté vient de leur intégration simultanée dans un objet compact et autonome, sous des contraintes qui s'opposent souvent.

La qualité audio demande une chaîne analogique à faible bruit. Or le Bluetooth place un émetteur radio juste à côté de cette chaîne sensible. Le format compact s'oppose ensuite à l'autonomie, car une batterie de grande capacité prend de la place. Le routage d'une carte multicouche doit également isoler les signaux audio analogiques des signaux numériques rapides et des plans d'alimentation. Le firmware doit enfin faire cohabiter des tâches aux exigences temporelles très différentes. Le flux audio ne tolère aucune interruption, alors que l'affichage et la lecture de la carte SD sont moins critiques.

Le problème central n'est donc pas le fonctionnement de chaque bloc pris isolément. Il s'agit de maîtriser les interfaces entre ces blocs. C'est à ces frontières, entre l'analogique et le numérique, entre le matériel et le logiciel, que se concentrent les compromis et les difficultés du projet.


\clearpage
\section{Objectifs}
Cette section définit les livrables nécessaires à la réalisation du projet. Les exigences fonctionnelles de l'appareil font l'objet du cahier des charges détaillé au chapitre~\ref{chap:analyse}.

L'objectif général consiste à concevoir intégralement un lecteur audio portable, compact et fonctionnel. La réalisation du projet suit les étapes suivantes :
\begin{itemize}
      \item établissement des spécifications techniques et de l'analyse fonctionnelle.
      \item conception de l'architecture matérielle, du schéma électrique et du routage d'un PCB multicouche.
      \item modélisation du boîtier mécanique.
      \item commande des composants et assemblage de la carte.
      \item développement du firmware assurant les fonctions requises.
      \item vérification et caractérisation de l'appareil terminé, du contrôle électrique de la carte aux mesures audio et d'autonomie.
      \item rédaction du présent rapport.
\end{itemize}


\section{Méthodologie}
La démarche suit un déroulement descendant, de la définition fonctionnelle vers la réalisation. Chaque décision technique découle ainsi d'un besoin établi en amont.

Le projet débute par une analyse fonctionnelle. Elle fixe ce que l'appareil doit faire, indépendamment de toute solution technique. Ces fonctions sont ensuite traduites en spécifications, qui guident le choix des composants et l'architecture matérielle. La conception matérielle précède le logiciel. Le schéma électrique, puis le routage du PCB, sont établis avant que le firmware ne puisse être testé sur la carte réelle. Ce séquencement impose d'anticiper, dès la conception de la carte, les contraintes que le logiciel rencontrera ensuite.

Le microcontrôleur retenu est l'ESP32. La carte est conçue sous Altium Designer, et le firmware sous ESP-IDF avec le système temps réel FreeRTOS. Ces trois choix ne sont pas posés d'avance. Ils sont établis et justifiés plus loin dans le rapport, une fois les spécifications fixées.

La vérification se déroule en deux temps. Pendant le développement, elle est incrémentale. Chaque bloc est testé seul, puis intégré au reste, jusqu'à la lecture audio continue et à la liaison Bluetooth. Sur l'appareil terminé, une campagne de mesure confronte chaque grandeur à un critère fixé avant l'essai. Ce critère provient du schéma, des fiches techniques ou du dimensionnement. Les performances de la sortie filaire sont relevées à l'analyseur audio, et l'autonomie par décharges complètes sur batterie.